\section{Bell-functional normalization and canonical quantity}\label{sec:normalization}
Let \(A_i=A_i^*=A_i^2\) and \(B_j=B_j^*=B_j^2\), for
\(i,j\in\{1,2,3\}\), be the event projectors of binary measurements. We use
\begin{equation}\label{eq:bell}
\begin{split}
\Bell={}&-A_2-B_1-2B_2
 +A_1B_1+A_1B_2+A_2B_1+A_2B_2\\
 &-A_1B_3+A_2B_3-A_3B_1+A_3B_2.
\end{split}
\end{equation}
Here \(A_iB_j\) means \(A_i\otimes B_j\); the omitted tensor factor in a
marginal is the identity. The convention uses event probabilities, not
\(\pm1\) observables.

\begin{definition}\label{def:St}
The number \(\St\) is the supremum of
\(\langle\psi,\Bell\psi\rangle\) over unit vectors
\(\psi\in\mathcal H_A\otimes\mathcal H_B\) and the six projectors above,
where both local Hilbert spaces are finite-dimensional and their dimensions
are not fixed in advance. Equivalently one may use normalized density
operators: a finite spectral decomposition expresses their values as convex
combinations of pure-state values for the same measurements.
\end{definition}

\begin{proposition}[Normalization identities]\label{prop:normalization}
The local bound of \eqref{eq:bell} is zero. In the event convention of
Mghirbi~\cite{Mghirbi2026paper}, the substitution
\[
 (P_1,P_2,P_3)=(A_2,A_1,A_3),\qquad
 (Q_1,Q_2,Q_3)=(B_2,B_1,B_3)
\]
identifies the Bell expression exactly, with scale one and shift zero.
Writing \(X_i=2P_i-I\), \(Y_j=2Q_j-I\), its involution version is
\begin{align}
 4\Bell&=\mathsf B-4I,\label{eq:involutions}\\
 \mathsf B&=X_1+X_2-Y_1-Y_2+X_1(Y_1+Y_2+Y_3)\notag\\
 &\quad+X_2(Y_1+Y_2-Y_3)+X_3(Y_1-Y_2).\notag
\end{align}
\end{proposition}
\begin{proof}
Substitution in
\(I_M=-P_1-2Q_1-Q_2+P_1Q_1+P_1Q_2+P_1Q_3+
P_2Q_1+P_2Q_2-P_2Q_3+P_3Q_1-P_3Q_2\)
gives every coefficient of \eqref{eq:bell}. Expanding
\(P_i=(I+X_i)/2\), \(Q_j=(I+Y_j)/2\) gives
\eqref{eq:involutions}. These maps are invertible on the stated projective
strategy class and preserve dimensions and the state. Exact evaluation on
the 64 deterministic assignments of six bits gives maximum zero and minimum
\(-3\); all-zero projectors attain zero. Convexity gives the same local bound
for mixtures. The coefficient identities and all 64 integer evaluations are
recomputed in the normalization certificate of Appendix~\ref{app:bounds}.
\end{proof}

We will use operators on \(\ell^2(\ZZ)\) as variational carriers. Their
appearance does not redefine \(\St\). The equality linking their variational
parameter to the finite tensor supremum is proved in
Proposition~\ref{prop:qstar}; it uses finite realizations and a specifically
restricted external upper-bound theorem.
